\section{Dataset-Native Data Preparation}
\label{sec:data}
\subsection{Choice of Primary Dataset}
The primary systems study uses TON-IoT~\cite{TONSource}. Its records combine IoT service telemetry, operating-system events, and networked interactions labeled with attack categories, making the dataset representative of the mixed-signal traffic that reaches an edge IDS at the service boundary rather than at the wire only. We deliberately restrict the physical dynamic-power study to a single dataset because coupling multiple encoded pipelines to one hardware campaign would confound feature dimensionality, class balance, and physical measurement; instead \cic{}~\cite{CICSource} and \nb{}~\cite{NBaIoTSource} are retained as supplementary detector and coverage audits and are not used to imply cross-dataset dynamic-power generalization.

\subsection{Frozen Preprocessing and Splits}
The TON-IoT pipeline is frozen and retains 32 native features. Imputation, scaling, and categorical encoding are fitted on training data only and applied deterministically to validation and test. The split groups unordered IP pairs so that traffic patterns are not divided across partitions; the grouping rule hashes ``$11\vert$'' concatenated with the sorted IP pair and thresholds the first sixteen SHA-256 hexadecimal digits modulo 100 (training below 60, validation 60--79, test 80--99). Training seed 11 fixes all model fits, so any single-seed conclusion is explicit rather than implied. Source-row identifiers and raw-field-string fingerprints are cross-partition disjoint, and an independent encoded-float32 audit finds \emph{zero} TON-IoT train--test vector matches, which addresses a commonly-cited leakage risk in IoT IDS benchmarks~\cite{Arp2022}.

\subsection{Auxiliary Dataset Audits}
The same ledger retains dataset-native \cic{} and \nb{} pipelines with 39 and 115 native features respectively. For \cic{}, seventeen attack labels remain training-only because their capture files are too few for a defensible test split; the actual test taxonomy is a binary attack-versus-normal partition over sixteen attack categories plus benign traffic. For \nb{}, physical device identifiers are the grouping unit, giving five training, two validation, and two test devices; a post-hoc byte-level audit records that 17.3\% of \nb{} test rows share exact float32 vectors with training rows (seven distinct vectors), reflecting the sensor-repetition structure of the dataset rather than a preprocessing fault. These coverage and encoded-overlap limitations are retained rather than used to imply cross-dataset dynamic-power generality; consequently the auxiliary datasets support only static classification and replay evidence.

\subsection{Class Balance and Replay Traces}
The TON-IoT test set contains a majority attack fraction, which is common in IoT benchmarks and which we do not rebalance because the physical study measures whole-device behavior on the natural distribution. To probe operating points under controlled prevalence shifts, the supplement retains twenty static replay traces per model (fifty per dataset before deduplication), with successive phase attack fractions $\{0.05, 0.35, 0.70, 0.40, 0.10\}$ over $\{100, 100, 120, 100, 80\}$ windows. These traces do not constitute a new detector-training exercise; they resample fixed held-out pools to expose how a frozen detector's F1 and balanced accuracy behave when prevalence, not distribution, changes.
